\documentclass{article}
\usepackage[T1]{fontenc}
\usepackage{lmodern}
\usepackage{graphicx}
\usepackage{amsmath,amssymb}
\usepackage[a4paper,margin=2cm]{geometry}
\usepackage[absolute,overlay]{textpos}
\setlength{\TPHorizModule}{1mm}
\setlength{\TPVertModule}{1mm}
\usepackage[indonesian]{babel}
\usepackage{hyperref}
\setlength{\emergencystretch}{2em}
% unit: Halaman 29

\begin{document}

% === Halaman 29 (python/native) ===

29

Q: What can a “bad guy” do?

A: a lot!

• eavesdrop: intercept messages

• actively insert messages into connection

• impersonation: can fake (spoof) source address in packet 

(or any field in packet)

• hijacking: “take over” ongoing connection by removing 

sender or receiver, inserting himself in place

• denial of service: prevent service from being used by 

others (e.g.,  by overloading resources)

Sumber: Chapter 8, Network Security

\end{document}
